\section{Future Work}
\label{sec:future}

\subsection{Scaling to extreme conversation lengths}
\label{sec:scale-future}

Our strategy comparison at 500K--10M tokens shows that frozen summary merges are not
triggered below $\sim$10M tokens ($\sim$168 cycles). Testing at 40M+ tokens would
force FrozenRanked into its hierarchical merge regime, revealing whether the
$\log_2(N)$ compression limit provides measurable benefit when merge pressure
is real --- or whether attention dilution dominates regardless.

\subsection{Importance-weighted compaction}
\label{sec:importance-weighted}

Score messages by importance before compaction, preserving high-value
exchanges (decisions, configurations) over low-value ones (acknowledgments).
Our category breakdown already shows that some fact types survive much
better than others; an explicit importance signal at compaction time should
amplify this effect on the parts of the conversation users actually care
about.

\subsection{RAG-augmented compaction}
\label{sec:rag}

Two-pass compaction: extract structured facts to a vector database + generate
a narrative summary. At query time, combine RAG retrieval with the summary.
Expected to address the ``keywords present but not retrievable'' gap.

\subsection{Lossless frozen with on-demand expansion}
\label{sec:lossless-frozen}

A natural variant of S3 Frozen: each frozen summary keeps a \textit{pointer} to the
verbatim text it summarises. The summary is what the model sees by default,
but a lightweight tool call (\code{expand(frozen\_id)}) can swap a frozen block
for its full original content when precise recall matters. Compared to pure
S3, this addresses the keyword-LLM gap exhibited in \S\ref{sec:dead-zone} without paying the
cost of carrying full context all the time. Compared to RAG (\S\ref{sec:rag}), it
preserves conversational structure (no arbitrary chunking) and uses the
model's existing summary as the index.

A more aggressive variant compacts continuously, one or a few turns at a
time, instead of waiting for a watermark threshold. This avoids the
``shock'' of large compaction passes (a major source of recall loss at 1M+
contexts, \S\ref{sec:strategy-results}) at the cost of more compaction calls. Existing systems like
Factory.ai's anchored iterative are in this regime; quantifying the
frequency-vs-granularity trade-off on the same evidence framework is a
natural extension of this paper.

\subsection{Graph-augmented context for million-token windows}
\label{sec:graph}

With 1M-token context windows now available (e.g.\ Claude Opus 4 with 1M
context), the cost of carrying an \textit{index} alongside summaries becomes
negligible --- 50 sessions $\times \sim$200 tokens of metadata is 1\% of context.
A ``pseudo-Graphiti in context'' approach attaches each frozen summary with
structured metadata (tags, time, actors, link to verbatim, link to other
related frozen blocks). The model can navigate the graph through its own
attention without making tool calls, enabling multi-hop reasoning over
preserved memory. This aligns with concurrent work on temporal knowledge
graphs for agent memory (Zep / Graphiti) but pushes the structure inside
the active context rather than into an external service.

\subsection{Implications of larger context windows}
\label{sec:large-windows}

The advent of 1M-token windows changes the operating point of compaction
but does not remove its relevance. With a 200K window, compaction triggers
roughly every 60K tokens ($\sim$30\% of the context) and is invoked frequently;
with 1M, it triggers less often but each invocation must compact $\sim$700K
tokens at once --- an order-of-magnitude larger and proportionally more
catastrophic when it goes wrong. The cost of \textit{holding} the larger context
also matters: KV cache memory grows linearly with window size, so even
when not running compaction, larger contexts have higher serving cost.
Future benchmarks should re-run our protocol at 1M windows to verify that
the strategy ordering and the variance phenomena documented here scale up
unchanged.

\subsection{Cross-model and cross-architecture validation}
\label{sec:cross-arch}

Our cross-model validation (\S\ref{sec:sonnet}) established that the findings generalize
from Claude Haiku 4.5 to Claude Sonnet 4.6 --- a model with substantially
different spatial recall characteristics (attenuated Lost-in-the-Middle,
$\sim$23pp range vs Haiku's $\sim$46pp). Testing on
non-Claude models (GPT-4, Gemini, open-source LLMs served locally) would
establish whether the findings hold across architectures. Preliminary
investigation suggests that local LLMs with shorter context windows
(32K--128K) could be tested with proportionally smaller contexts.
